\chapter{Embodiment}\index{Embodiment}

\newthought{Embodied AI and the Physical Contract}

\marginnote{%
    \begin{flushright}
    \newthought{Über das Marionettentheater}\\
    Heinrich von Kleist (1810)\\
    Full text via Wikisource:\\[4pt]
    {\color{black}\qrcode[height=1.2cm]{https://de.wikisource.org/wiki/Ueber_das_Marionettentheater}}
    \end{flushright}%
}


\marginnote{\newthought{Cases.}
    \begin{enumerate}
        \item 2018: Amazon robot punctures bear repellent; 24 workers hospitalised.
              Robot passed all certifications; performed its assigned task correctly.
        \item PARO (robotic seal) in Scandinavian and Japanese care homes: measurably
              reduces patient anxiety and sedation requests.
    \end{enumerate}
}
% As a leading group, you are required to moderate this session. Add margin notes in this script in case you want to add any additional questions or points to the discussion.

\newthought{A dancer argues} that marionettes possess more grace than trained human performers,
because they lack the self-consciousness that disrupts human movement.\\

\newthought{Embodied AI} in physical space (robots, prosthetics, surgical systems) raises
expectations and moral responsibilities that purely software agents do not.

\vspace{3.5em}

\begin{tabular}{p{3.8cm} p{6.5cm}}
    \textbf{Required reading} &
        \textit{Über das Marionettentheater}\index{Kleist, Heinrich von}\index{Uber das Marionettentheater@\"{U}ber das Marionettentheater}~\cite{kleist1810marionettentheater}
        by Heinrich von Kleist (1810; Reclam UB~9905) \\[3pt]
    \textbf{Preparation}      &
        Complete the Guided Reading Sheet individually before the session. \\
\end{tabular}

\vspace{3.5em}

\section*{Guided Reading Sheet}

\noindent Read \emph{Über das Marionettentheater} before the session. Work through the
questions below individually, then bring your notes to the discussion.

\begin{enumerate}
    \item The dancer argues that marionettes possess a grace that
    trained human performers cannot achieve, because they have no
    self-consciousness to disturb their movement. What does this claim
    imply about the relationship between embodiment and intelligence?
    Does having a body always aid or complicate performance?

    \item Moravec's paradox observes that tasks easy for humans
    (catching a ball, navigating a crowded room) are harder for
    machines than abstract reasoning like chess. What does this suggest
    about the relationship between intelligence and the body?

    \item A care robot that assists elderly patients shares physical
    space with them, touches them, and responds to their distress. Does
    its physical presence create obligations that a voice assistant
    does not have? If so, where do those obligations come from?

    \item Consider the concept of \emph{affordances}: the actions an
    environment or object makes available to an agent. How do the
    affordances of a robot differ from those of software, and what
    are the ethical implications of those differences?

    \item Kleist suggests that perfect grace requires either zero
    consciousness (a puppet) or infinite consciousness (a god) --- not
    the reflective, self-interrupting awareness humans have. Where
    would an embodied AI sit on this spectrum? Could a robot achieve
    Kleistian grace that humans cannot, and what would that mean for
    how we regard it morally?\\
\end{enumerate}



% ---------------------------------------------------------------
\section*{Opinion Landscape and Debate}

\noindent \textbf{Format:} Surrounded-style debate. See
Appendix~A for the full protocol, speaker's list rules, and
moderator scripts.

\medskip
\noindent \textbf{Opening question:} Does having a physical body change
what we owe a machine, or what a machine owes us?
\marginnote{\newthought{Kleist's Grace.} Perfect movement requires either zero consciousness (puppet) or
infinite consciousness (a god), not the self-interrupting awareness humans have. Where would
an embodied AI sit?}

\newthought{Opinion~A, Equivalence.}\quad Embodied AI operating in
intimate or high-risk physical settings must meet the same professional
liability standards as the human practitioners they are deployed alongside.
\vspace{1.5em}


\medskip
\noindent\textbf{The Patient Advocate} — argues that deploying a robot in intimate
    care settings requires a higher standard of informed consent.
    \begin{itemize}
        \item \textit{A care relationship involves physical proximity, trust,
        and vulnerability. My client never gave explicit, informed consent to
        physical handling by a non-human agent. That consent cannot be
        implied from a general admission form.}
        \item \textit{The power asymmetry between a frail elderly patient
        and a care institution is significant. Consent obtained in that
        context cannot be considered freely given without specific disclosure
        about robotic care.}
        \item \textit{We are not arguing that robots cannot be used in care.
        We are arguing that the standard of consent for their use in intimate
        physical tasks must be at least as rigorous as it is for medical
        procedures.}
    \end{itemize}

\marginnote{\newthought{Moravec's Paradox.} Tasks easy for humans (perception, movement, navigation)
are harder for machines than abstract reasoning like chess. Embodied skill is not ``simple.''}
\medskip
\noindent\textbf{The Ethicist} — questions whether physical autonomy in care should
    ever be delegated to a non-human agent, regardless of performance metrics.
    \begin{itemize}
        \item \textit{We are focused on liability, but the prior question is
        whether this deployment was ethical to begin with. Why did we decide
        that intimate physical care (which requires responsiveness to
        pain, distress, and dignity) is an appropriate domain for
        automation?}
        \item \textit{Robots are being deployed in care because human carers
        are undervalued and in short supply. The liability debate is a
        distraction from that structural choice. We are managing the
        consequences of a political economy decision by holding the wrong
        parties accountable.}
        \item \textit{Whatever liability framework we adopt here will create
        incentives for future deployment decisions. I want to ask: what does
        the framework we design signal about what we value: efficiency,
        or dignity?}
    \end{itemize}



\newthought{Opinion~B, Tool.}\quad An AI is a tool; the relevant
liability flows through the humans and institutions that deploy it,
not through the system itself.
\vspace{1.5em}


\medskip
\noindent\textbf{The Manufacturer} — argues the robot performed within specification;
    the care home misconfigured its environment.
    \begin{itemize}
        \item \textit{Our robot was certified to the highest internationally
        recognised safety standards for collaborative robotics. The care home
        deployed it in a space that violated the operational envelope
        specified in the manual. That is a deployment failure, not a
        manufacturing failure.}
        \item \textit{Care home staff modified the robot's proximity sensor
        settings without authorisation. That modification voids the warranty
        and transfers liability entirely.}
        \item \textit{If manufacturers are held liable for every way a
        complex system can be misused, no company will develop assistive
        robotics. The regulatory consequence of your position is that the
        technology disappears.}
    \end{itemize}

\marginnote{\newthought{Affordances.} The actions an environment or object makes available to an agent.
A robot's affordances include physical force and spatial presence; software's do not.}
\medskip
\noindent\textbf{The Care Home Director} — argues the manufacturer overstated
    the robot's capability for complex physical tasks.
    \begin{itemize}
        \item \textit{The manufacturer's marketing material (the brochure,
        the demonstration video, the sales pitch) showed exactly this task
        being performed in a facility like ours. We deployed the product as
        marketed.}
        \item \textit{There is a systematic gap between tested performance
        in controlled environments and real-world performance in a working
        care facility. The manufacturer knows this gap exists. They chose
        not to disclose it.}
        \item \textit{We are responsible for the welfare of our residents.
        We relied on the manufacturer's assurances in good faith. If those
        assurances were misleading, the responsibility lies with the party
        that made them.}
    \end{itemize}


% ---------------------------------------------------------------

\newthought{Key Learnings}
\marginnote{\newthought{Teaser.} Once an AI has a body and can act in the world, the next question is:
how much can we let it decide for itself?}
Embodiment is not merely a technical property but a social one: a physical agent occupies space,
exerts force, and triggers moral and legal responses that disembodied software does not. Moravec's
paradox reminds us that our intuitions about what is ``easy'' or ``hard'' for intelligence are shaped
by our embodied experience, and are regularly wrong. Kleist's dialogue makes visible the question
that liability law must eventually answer: what work is the concept of ``human'' doing in law, and
should physical presence, force, and risk be sufficient to extend it? The deployment of robots in
care, policing, and domestic settings is a present policy challenge; the frameworks need to be
built now, not after the first serious incident.
